% !Mode:: "TeX:UTF-8"
% 第 2 章，照范例第 4 章节选。公式按章编号，引用写 \eqref；公式下面的“式中”说明用
% eqdenote 环境，一条一个 \item，符号与说明之间写 &，破折号与对齐由模板排。

\chapter{基于FLUENT软件的轴承静态特性研究}[Research on static characteristics of bearing based on FLUENT software]

\section{引言}[Introduction]

利用现成的商用软件来研究流场，可以免去对N-S方程求解程序的……，本章在此基础上用\abbr{FEM}求解局部多孔质区域的流场，……

\section{边界条件的设定}[Initialization of boundary conditions]

本文采用……，则每一个方向上的……由公式\eqref{eq:2-1}、\eqref{eq:2-2}求得：
\begin{equation}\label{eq:2-1}
  \alpha = \frac{D_{\mathrm{p}}^2}{150}\frac{\psi^3}{(1-\psi)^2}
\end{equation}
\begin{equation}\label{eq:2-2}
  C_2 = \frac{3.5}{D_{\mathrm{p}}}\frac{(1-\psi)}{\psi^3}
\end{equation}
\begin{eqdenote}
  \item $D_{\mathrm{p}}$ & 多孔质材料的平均粒子直径（m）；
  \item $\psi$ & 孔隙度（孔隙体积占总体积的百分比）；
  \item $C_2$ & 特征渗透性或固有渗透性（m$^2$）。
\end{eqdenote}

……

\section{本章小结}[Brief summary]

……
